\Askhsh{36226}{2}
\begin{ylh}{1}
    \item 3.7. Κύκλος - Μεσοκάθετος - Διχοτόμος.
\end{ylh}
% ===================================================================
%  Αρχείο: 36226.tex (Fragment)
%  Αυτόματη μετατροπή μέσω doc2latex.py
% ===================================================================

ΘΕΜΑ 3

Στο παρακάτω σχήμα έχουμε το χάρτη μίας περιοχής όπου είναι κρυμμένος ένας θησαυρός. Οι ημιευθείες Αx και Αy παριστάνουν δύο ποτάμια και στα σημεία Β και Γ βρίσκονται δύο πλατάνια. Ο πλάτανος που βρίσκεται στο σημείο Β έχει μικρότερη απόσταση από το σημείο Α, σε σχέση με την απόσταση που έχει από το σημείο Α ο πλάτανος που βρίσκεται στο σημείο Γ.

Να προσδιορίσετε γεωμετρικά τις δυνατές θέσεις του θησαυρού, αν είναι γνωστό ότι:

\begin{alist}
\item ο θησαυρός ισαπέχει από τα δύο πλατάνια. \monades{9}

\item ο θησαυρός ισαπέχει από τα δύο ποτάμια. \monades{9}

\item ο θησαυρός ισαπέχει από τα δύο πλατάνια και ισαπέχει και από τα δύο ποτάμια.

\monades{7}

Να αιτιολογήσετε την απάντησή σας σε κάθε περίπτωση.


\begin{center}
\begin{tikzpicture}[scale=1]
\tkzDefPoint(0,0){A}
\tkzDefPoint(5,0){x_axis}
\tkzDefPointBy[rotation=center A angle 50](x_axis) \tkzGetPoint{y_axis}
\tkzDefPoint(2,0){B}
\tkzDefPointBy[rotation=center A angle 50](x_axis) \tkzGetPoint{C_dir}
\tkzDefPointBy[homothety=center A ratio 0.8](C_dir) \tkzGetPoint{C}

\draw[l3, ->] (A)--(x_axis) node[right]{$x$};
\draw[l3, ->] (A)--(y_axis) node[above]{$y$};
\tkzDrawPoints(B,C)
\tkzLabelPoint[below](B){$B$}
\tkzLabelPoint[above left](C){$\Gamma$}

\node at (1,-1) {Ποτάμι $Ax$};
\node at (-1,2) {Ποτάμι $Ay$};
\end{tikzpicture}
\end{center}
\end{alist}
